\section{Negative recursive attempts: V16--V18}

A central reason to publish the full sequence rather than only V19--V20 is that the failed recursive attempts expose the failure modes that the successful protocol eventually addresses.

\subsection{V16 and V17: bounded policy-class saturation}

After the G0$\rightarrow$G1 L2 result, V16 attempts to author G2 within the inherited bounded policy family. It fails to produce a successor that passes the required gate. V17 then increases chronological replay resolution by rebuilding 13 prefix worlds from the real archive. The result remains negative. The evidence register concludes that the current bounded meta-policy family is saturated on the admitted real tree.

This is not treated as almost G2. It is retained as evidence that a narrow parameter neighborhood can exhaust its useful variation even when the broader recursive question remains open.

\subsection{V18: executable-policy expansion and an overfitting failure}

V18 expands the editable object from the bounded parameter vector to executable policy code, subject to an AST/ABI sandbox. Four independent isolated proposer attempts all discover the same public strategy: on the 13 real prefix worlds, stop immediately at the already-observed high-quality plateau. Each child improves mean public replay core from 943 to 953 and reduces represented requests from 1 to 0. All four are decision-equivalent despite different source hashes.

The selected source is frozen before a fresh 24-world holdout is allocated. The result is completely neutral:
\[
0\ \text{wins},\quad 0\ \text{losses},\quad 24\ \text{ties}.
\]
Parent and child both obtain mean core 666.125 and mean represented requests 5.0. V18 therefore fails.

The failure is informative. A single real replay regime made a plateau-specific shortcut look universally useful. The external holdout prevents that shortcut from being admitted as G2. V19 is designed only after this negative is public.

\subsection{Why V18 matters for interpretation}

RRSI independently highlights a closely related concern: agent-harness evolution can show large in-distribution improvements that shrink outside the evolve set \cite{xia2026rrsi}. V18 is a concrete instance of that risk in Genesis. The project does not claim that V19 eliminates overfitting in general; instead, V19 expands the development distribution and reserves a new untouched holdout.
